%
%------------------------------------------------------------------------------------
%	AUTORI: Morati Mirko, Noè Murr
%------------------------------------------------------------------------------------
%

%------------------------------------------------------------------------------------
%	PREAMBOLO
%------------------------------------------------------------------------------------
\documentclass[a4paper,11pt]{article}
\usepackage[T1]{fontenc}
\usepackage[utf8x]{inputenc}
\usepackage[italian]{babel}
\usepackage{amsmath}
\usepackage{calc}
\usepackage[a4paper]{geometry}
\usepackage[usenames,dvipsnames]{xcolor}
\usepackage{fancyhdr} % Required for custom headers
\usepackage{lastpage} % Required to determine the last page for the footer
\usepackage{extramarks} % Required for headers and footers
%\usepackage[usenames,dvipsnames]{color} % Required for custom colors
\usepackage{graphicx} % Required to insert images
\usepackage{listings} % Required for insertion of code
\usepackage{courier} % Required for the courier font
\usepackage{lipsum} % Used for inserting dummy 'Lorem ipsum' text into the template
\usepackage{pdfpages}
\usepackage{listings}
\usepackage{mdframed}
\usepackage{enumitem}
%\usepackage{minted}
\usepackage{float}
\usepackage{array}
\usepackage{zref-xr}
\usepackage{zi4}
\usepackage{tikz}
\usetikzlibrary{automata, arrows, positioning, calc, matrix, shapes.geometric, chains}
\usepackage{verbatim}
\usepackage{color}   %May be necessary if you want to color links
\usepackage{hyperref}
\hypersetup{
	colorlinks,
	citecolor=black,
	filecolor=black,
	linkcolor=black,
	urlcolor=black
}


\newenvironment{framescelte}[1]
  {\mdfsetup{
    frametitle={\colorbox{white}{\space#1\space}},
    innertopmargin=6pt,
    frametitleaboveskip=-\ht\strutbox,
    frametitlealignment=\center
    }
  \begin{mdframed}%
  }
  {\end{mdframed}}


% Margins
\topmargin=-0.45in
\evensidemargin=0in
\oddsidemargin=0in
\textwidth=6.25in
\textheight=9in
\headsep=0.25in

\linespread{1.1} % Line spacing


% Set up the header and footer
\pagestyle{fancy}
%\lhead{\hmwkAuthorName} % Top left header
\chead{\hmwkTitle} % Top center head
\rhead{\firstxmark} % Top right header
\lfoot{\lastxmark} % Bottom left footer
%\cfoot{} % Bottom center footer
%\rfoot{Page\ \thepage\ of\ \protect\pageref{LastPage}} % Bottom right footer
\renewcommand\headrulewidth{0.4pt} % Size of the header rule
\renewcommand\footrulewidth{0.4pt} % Size of the footer rule

%------------------------------------------------------------------------------------
%	TITOLO
%------------------------------------------------------------------------------------

\newcommand{\hmwkTitle}{Elaborato\ Assembly}
\newcommand{\hmwkClass}{Architettura degli Elaboratori}
\newcommand{\hmwkAuthorName}{Mirko Morati,\ Noè Murr}

%------------------------------------------------------------------------------------
%	TITLE PAGE
%------------------------------------------------------------------------------------

\title{
\vspace{2in}
\textmd{\textbf{\hmwkClass:\\ \hmwkTitle}}\\
\vspace{0.1in}\large{\textit{\hmwkAuthorName}}
\vspace{3in}
}

%\author{}

%------------------------------------------------------------------------------------


%------------------------------------------------------------------------------------
%	FORMATTAZIONE ELEMENTI
%------------------------------------------------------------------------------------

\renewcommand{\labelitemi}{$\cdot$}

\newcommand{\Assembly}{\texttt{Assembly} }

\newcommand{\itemtt}[1]{\item \texttt{#1}}

\newcommand{\myparagraph}[2]{
	\begin{table}[!ht]
		\begin{tabular}{p{0.175\linewidth} | p{0.8\linewidth}}
			\texttt{#1} & #2
		\end{tabular}
	\end{table}
			}
			

%------------------------------------------------------------------------------------
%	TIKZ DEFINIZIONI
%------------------------------------------------------------------------------------
\tikzstyle{base} = [text width = 10em]
\tikzstyle{startstop} = [rectangle, rounded corners, base, text centered, draw=black]
\tikzstyle{process} = [rectangle, base, text centered, draw=black]
\tikzstyle{decision} = [
diamond,
aspect = 2,
draw,
text width=8em,
text badly centered,
inner sep=0pt]
\tikzstyle{io} = [trapezium, base, trapezium left angle=70, trapezium right angle=110, text centered, draw=black]

\tikzstyle{arrow} = [draw, thick,->,>=stealth]

%------------------------------------------------------------------------------------
%	LISTINGS DEFINIZIONI
%------------------------------------------------------------------------------------

\definecolor{mygreen}{rgb}{0,0.6,0}
\definecolor{mygray}{rgb}{0.5,0.5,0.5}
\definecolor{mymauve}{rgb}{0.58,0,0.82}


\lstdefinelanguage{MyAssembler}{
	morecomment=[l][\color{mygray}]{\#},
	morekeywords=[2]{eax, ebx, ecx, edx, edi, esp, ebp, al, bl, cl, dl, ah, bh, ch, dh},
	%	morekeywords=[3]{\$, 0, 1, 2, 3, 4, 5, 6, 7, 8, 9},
	morekeywords=[4]{popl,popb,pushl,pushb,esi,jne,jn,call,int,movl,movb,xorl,xorb,subl,subb,addl,addb,mull,mulb,incl,cmpb,cmpl,jl,jg,jge,jle,je,jmp,ret,decl,leal},
	morekeywords=[5]{section,data,type,long,bss,text,globl,equ,ascii,asciz,include},
	morestring=[b][\color{mymauve}]{"},
	morekeywords=[6][\color{Violet}]{},
	morekeywords=[7]{input\_fd,output\_fd,init,reset,rpm,alm,numb,mod}
}

\lstdefinestyle{MyAsm}{ %
	backgroundcolor=\color{white},   % choose the background color; you must add \usepackage{color} or \usepackage{xcolor}
	basicstyle=\footnotesize\ttfamily,        % the size of the fonts that are used for the code
	breakatwhitespace=false,         % sets if automatic breaks should only happen at whitespace
	breaklines=true,                 % sets automatic line breaking
	captionpos=b,                    % sets the caption-position to bottom
	commentstyle=\color{mygray},    % comment style
	deletekeywords={...},            % if you want to delete keywords from the given language
	escapeinside={\%*}{*)},          % if you want to add LaTeX within your code
	extendedchars=true,              % lets you use non-ASCII characters; for 8-bits encodings only, does not work with UTF-8
	literate={è}{{\'e}}1 {ì}{{\'i}}1 {é}{{\'e}}1  {ò}{{\'o}}1 {à}{{\'a}}1,
	frame=single,	                   % adds a frame around the code
	keepspaces=true,                 % keeps spaces in text, useful for keeping indentation of code (possibly needs columns=flexible)
	keywordstyle=\color{blue},       % keyword style
	keywordstyle=[2]\color{BrickRed},
	%	keywordstyle=[3]\color{green},
	keywordstyle=[4]\color{blue},	
	keywordstyle=[5]\color{Green},
	keywordstyle=[7]\color{Orange},
	%language={[x86masm]Assembler},                 % the language of the code
	otherkeywords={...},           % if you want to add more keywords to the set
	numbers=left,                    % where to put the line-numbers; possible values are (none, left, right)
	numbersep=5pt,                   % how far the line-numbers are from the code
	numberstyle=\tiny\color{mygray}, % the style that is used for the line-numbers
	rulecolor=\color{black},         % if not set, the frame-color may be changed on line-breaks within not-black text (e.g. comments (green here))
	showspaces=false,                % show spaces everywhere adding particular underscores; it overrides 'showstringspaces'
	showstringspaces=false,          % underline spaces within strings only
	showtabs=false,                  % show tabs within strings adding particular underscores
	stepnumber=1,                    % the step between two line-numbers. If it's 1, each line will be numbered
	stringstyle=\color{mymauve},     % string literal style
	tabsize=4,	                   % sets default tabsize to 2 spaces
	escapeinside={(*@}{@*)},
	%title=\lstname                   % show the filename of files included with \lstinputlisting; also try caption instead of title
}



\begin{document}
	\clearpage
%	\maketitle
	\begin{titlepage}
		\centering
		\vspace*{\fill}
		\includegraphics[width=0.15\textwidth]{logo.png}\par\vspace{1cm}
		{\scshape\LARGE Università degli Studi di Verona \par}
		\vspace{1cm}
		{\scshape\Large Architettura degli Elaboratori\par}
		\vspace{1.5cm}
		\\ \line(1,0){250} \\
		{\huge\bfseries Elaborato Assembly\par}
		\\ \line(1,0){250} \\
		\vspace{2cm}
		{\Large\itshape Mirko Morati, Noè Murr\par}
		\vspace{5cm}
		\vspace*{\fill}
		% Bottom of the page
		{\large \today\par}
	\end{titlepage}
	\thispagestyle{empty}
	\newpage
	\tableofcontents
	\newpage
	
	\section{Premessa}
	Per la corretta lettura del pdf si informa che sono stati inseriti dei riferimenti per ogni etichetta e variabile trattata in modo da poter vedere e approfondire il corrispettivo codice nel programma. 
	
	\section{Descrizione del progetto}
	Si vuole realizzare un programma \Assembly per il monitoraggio di un motore a combustione interna il quale, ricevuto come ingresso il numero di giri/minuto del motore, fornisca in uscita la modalità di funzionamento corrente del motore: \textit{Sotto Giri, Ottimale, Fuori Giri}. Il programma deve contare e visualizzare in uscita il numero dei secondi trascorsi nella modalità di funzionamento attuale ed inoltre attivare il segnale di allarme nel caso in cui il motore si trovi in modalità \textit{Fuori Giri} da più di 15 secondi.
	
	\vspace {5mm}
	Di seguito verranno descritte le funzioni presenti in ogni file del programma, etichette, eventuali variabili e loro scopo. 
	
	\section{syscall.inc}
	Header file contenente la definizione di alcune costanti, tramite la pseudo-operazione \texttt{.equ}, relative alle chiamate di sistema e ad alcuni standard utilizzati in tutti i file e riportati di seguito:
	\begin{table}[!h]
		\begin{tabular}{| >{\ttfamily}l | c |}
			\hline
			\hyperref[s:1]{SYS\_EXIT} & 1 \\ \hline
			SYS\_READ & 3 \\ \hline
			SYS\_WRITE & 4 \\ \hline
			SYS\_OPEN & 5 \\ \hline
			SYS\_CLOSE & 6 \\ \hline
			STDIN & 0 \\ \hline
			STDOUT & 1 \\ \hline
			STDERR & 2			 \\ \hline
			SYSCALL & 0x80 \\ \hline
		\end{tabular}
	\end{table}
	
	\begin{framescelte}{Scelte Progettuali}
		Nonostante il codice \Assembly sia fortemente legato alla macchina sottostante, si è deciso di utilizzare un "trucchetto" per permetterne, per quanto possibile, la portabilità agevole su un'altra piattaforma. Infatti sarà sufficiente cambiare i codici delle \texttt{SYS\_CALL} nel file \texttt{syscall.inc} per utilizzare le \texttt{SYS\_CALL} di un altro sistema operativo. Questo vale solo su sistemi unix-like poiché i parametri delle funzioni essenziali (read, write, \ldots) sono comuni grazie allo standard \texttt{POSIX}.
	\end{framescelte}
	
	\newpage
	
	\section{main.s}
	File principale del programma. 
	
	\subsection{Flowchart}
	\begin{center}
		
		\begin{tikzpicture}[node distance = 2cm, auto]
		
		
		\node	(start)	   [startstop] 					{Inizio};
		\node	(i1) 		[io, below of = start] 				 {eax $\leftarrow$ \#argomenti};
		\node	(d1) 	  [decision, below of = i1] 	  {eax == 3 ?};
		\node	(e1)    [startstop, right of = d1, node distance = 6cm] 	 {Stampa messaggio ed esci con errore 1};
		\node	(i2)	[io, below of = d1, node distance = 2.5cm]		{input\_file\_name output\_file\_name};
		\node	(p1)  [process, below of = i2] 	  {call \_open\_files};
		\node	(p2)	[process, below of = p1]	{call \_read\_line};
		\node	(d2)   [decision, below of = p2]	  {ebx == -1 ?};
		\node   (p3)	[process, right of = d2, node distance = 6cm]		{call \_close\_files};
		\node	(e0)   [startstop, below of = p3]		{Esci con codice 0};
		\node	(p4)	[process, below of = d2, node distance = 2.5cm]		{call \_check};
		\node 	(p5)	[process, below of = p4]		{call \_write\_line};
		
		\path	[arrow]		(start)	--	(i1);
		\path	[arrow]		(i1)	--	(d1);
		\path	[arrow]		(d1)	--	node {Si} (i2);
		\path	[arrow]		(d1)	--	node {No} (e1);
		\path	[arrow]		(i2)	--	(p1);
		\path	[arrow]		(p1)	--	(p2);
		\path	[arrow]		(p2)	--	(d2);
		\path	[arrow]		(d2)	--	node {Si} (p4);
		\path	[arrow]		(d2)	--	node {No} (p3);
		\path	[arrow]		(p3)	--	(e0);
		\path	[arrow]		(p4)	--	(p5);
		\path	[arrow]		(p5.west) node [anchor =south west] {} --+(-2cm,0)  |-(p2);
		
		\end{tikzpicture}
		
	\end{center}
	
	\newpage
	
	\subsection{Variabili Globali}
	\begin{itemize}
		\itemtt{\hyperref[v:1:1]{input\_fd}}: Contiene il descrittore del file di input;
		\itemtt{\hyperref[v:1:2]{output\_fd}}: Contiene il descrittore del file di output;
		\itemtt{\hyperref[v:1:3]{init}}: Contiene il valore del segnale INIT corrente;
		\itemtt{\hyperref[v:1:4]{reset}}: Contiene il valore del segnale RESET corrente;
		\itemtt{\hyperref[v:1:5]{rpm}}: Contiene il valore del segnale RPM corrente;
		\itemtt{\hyperref[v:1:6]{alm}}: Contiene il valore del segnale ALM corrente;
		\itemtt{\hyperref[v:1:7]{mod}}: Contiene il valore del segnale MOD corrente;
		\itemtt{\hyperref[v:1:8]{numb}}: Contiene il valore del segnale NUMB corrente.
	\end{itemize}
	
	\begin{framescelte}{Scelte Progettuali}
		Si è scelto di utilizzare in modo estensivo le variabili globali in quanto in un progetto di così ridotte dimensioni si aumenta notevolmente la leggibilità del codice e nel contempo se ne riduce la complessità di scrittura. Se il progetto avesse richiesto computazioni più complesse questo approccio non sarebbe stato l'ideale, poiché sarebbero state compromesse le prestazioni totali dell'applicazione.
	\end{framescelte}
	\subsection{Variabili Locali}
	\begin{itemize}
		\itemtt{\hyperref[v:1:9]{usage}}: Stringa per la descrizione del corretto utilizzo del programma;
		\itemtt{\hyperref[v:1:10]{USAGE\_LENGTH}}: Contiene la lunghezza della stringa, necessaria per la stampa.
	\end{itemize}
	
	\subsection{Funzioni ed Etichette}
		\begin{itemize}
		\itemtt{\hyperref[e:1:1]{\_start}}: Punto di entrata del programma. Si occupa di controllare che il numero di parametri sia corretto, in caso contrario stampa la stringa \texttt{usage} e termina. Dopo il controllo chiama la funzione \texttt{\_open\_files} definita nel file \texttt{open\_files.s}. 
		\itemtt{\hyperref[e:1:2]{\_main\_loop}}: Loop principale. Viene chiamata la funzione \texttt{\_read\_line} definita nel file \texttt{read\_line.s}, nel caso in cui il contenuto del registro \texttt{ebx} sia equivalente a -1 significa che il file di input è terminato (\textbf{EOF} $\rightarrow$ End Of File) quindi salta a \texttt{\_end}, altrimenti chiama la funzione \texttt{\_check} definita nel file \texttt{check.s} e la funzione \texttt{\_write\_line} definita nel file \texttt{write\_line.s}, dopodiché riesegue il ciclo.
		\itemtt{\hyperref[e:1:3]{\_end}}: Si occupa di chiudere tutti i file aperti e della corretta uscita dal programma tramite la chiamata di sistema \texttt{exit}.
		\itemtt{\hyperref[e:1:4]{\_show\_usage}}: Nel caso in cui i parametri non siano corretti stampa a video la stringa \texttt{usage} e termina il programma segnalando errore con il codice 1.
	\end{itemize}  
	
	\section{open\_files.s}
	Contiene la funzione che si occupa di aprire i file in modo corretto.
	
	\subsection{Variabili Locali} 
	\begin{itemize}
		\itemtt{\hyperref[v:2:1]{error\_opening\_files}}: Stringa di errore in caso di errata apertura dei file (e.g. file mancante, file corrotto, \ldots).
		\itemtt{\hyperref[v:2:2]{ERROR\_OPENING\_LENGTH}}: Costante che contiene la lunghezza della stringa di errore.
	\end{itemize}
	
	\subsection{Funzioni ed Etichette}
	\begin{itemize}
		\itemtt{\hyperref[e:2:1]{\_open\_files}}: Si occupa di aprire i file e gestisce eventuali errori. In caso il file di output non esistesse, questo viene creato in automatico con permessi di lettura e scrittura. I descrittori ottenuti vengono salvati nelle corrispondenti variabili globali. 
		\itemtt{\hyperref[e:2:2]{\_error\_opening\_files}}: In caso di errore viene stampato su \texttt{STDERR} la stringa opportuna, dopodiché il programma viene terminato con codice di errore 2. 
	\end{itemize}
	
	\section{read\_line.s}
	Contiene la funzione che si occupa di leggere ed interpretare una riga per volta del file di input.
	\subsection{Variabili Locali}
	\begin{itemize}
		\itemtt{\hyperref[v:3:1]{input\_buff}}: Buffer di dimensione \texttt{INPUT\_BUFF\_LEN} che conterrà i caratteri della riga letta dal file di input.
		\itemtt{\hyperref[v:3:2]{INPUT\_BUFF\_LEN}}: Dimensione del buffer.
	\end{itemize}
	
	\begin{framescelte}{Scelte Progettuali}
		Si è scelto di utilizzare un buffer di dimensione 9 byte corrispondente alla lunghezza di una singola riga. Questa scelta è molto importante perché permette di non realizzare soluzioni hardcoded. La consegna specifica che non verranno utilizzati file di input con più di 100 righe: sfruttando tale informazione si sarebbe potuto scrivere un codice semplicistico che leggesse in un buffer l'intero file. Soluzione poco elegante e mal ottimizzata. Un'altra strada prendibile sarebbe stata quella di leggere un byte alla volta dal file di input, ma questa avrebbe richiesto un numero spropositato e inutile di chiamate al kernel. Con la lettura di una riga per volta il buffer rimane snello e comodo da manipolare. Questa soluzione permette all'applicazione di lavorare con file di dimensione arbitraria, inoltre lo riteniamo un metodo ottimale in quanto stabilisce un compromesso tra l'uso di \texttt{SYS\_CALL} e flessibilità del codice. Oltre a questo si sarebbe potuto utilizzare un ulteriore metodo: quello di mappare in memoria il file ed accedervi come in un array, ma questo sforava dalle nozioni fornite dal corso.
	\end{framescelte}
	
	\subsection{Funzioni ed Etichette}
	\begin{itemize}
		\itemtt{\hyperref[e:3:1]{\_read\_line}}: Legge dal file una riga tramite la chiamata di sistema \texttt{read} e si occupa di richiedere la traduzione dei caratteri letti in interi mediante la funzione \texttt{\_atoi} definita nel file \texttt{atoi.s} salvandoli successivamente nelle rispettive variabili globali. In caso i caratteri letti siano pari a 0 salta all'etichetta \texttt{\_eof};
		\itemtt{\hyperref[e:3:2]{\_eof}}: In caso di \texttt{EOF} mette -1 in \texttt{ebx} e ritorna.
	\end{itemize}
	
	
	\section{atoi.s}
	Contiene la funzione che si occupa di convertire una serie di caratteri \texttt{ASCII} in un numero intero. 
	\subsection{Funzioni ed Etichette}
	\begin{itemize}
		\itemtt{\hyperref[e:4:1]{\_atoi}}: Vengono inizializzati i registri necessari alla conversione;
		\itemtt{\hyperref[e:4:2]{\_atoi\_loop}}: Loop principale, converte la stringa puntata dal registro \texttt{edi} in un intero salvato e ritornato in \texttt{eax}. 
	\end{itemize}
	
	\section{check.s}
	Contiene la funzione che si occupa di settare sulla base dei valori di input e dei valori del ciclo precedente i corretti parametri delle variabili \texttt{alm, mod, numb}. 
	\subsection{Funzioni ed Etichette}
	\begin{itemize}
		\itemtt{\hyperref[e:5:1]{\_check}}: In base ai valori di \texttt{init} e \texttt{rpm} si occupa di saltare all'etichetta corretta. 
		\itemtt{\hyperref[e:5:2]{\_fg / \_sg / \_opt}}: Etichette corrispondenti alle modalità di funzionamento previste dalle specifiche. Si occupano di settare i corretti valori di \texttt{alm, mod, numb}. L'unica modalità che necessita di una gestione particolare è \texttt{fg} in cui bisogna settare l'eventuale allarme. 
		\itemtt{\hyperref[e:5:3]{\_reset\_numb}}: Nel caso in cui sia stata cambiata la modalità di funzionamento oppure il valore della variabile \texttt{reset} sia pari a 1, viene settata la modalità corretta, resettato il conteggio \texttt{numb} e "spento" \texttt{alm}.
		\itemtt{\hyperref[e:5:4]{\_set\_alm}}: Se il motore è nella modalità \texttt{fg} da più di 15 secondi, viene "acceso" l'allarme portando il valore di \texttt{alm} a 1.
		\itemtt{\hyperref[e:5:5]{\_init\_0}}: Se il valore di \texttt{init} è pari a 0 tutte le variabili di output vengono poste a 0 dal momento che il motore è spento. 
		\itemtt{\hyperref[e:5:6]{\_end\_check}}: Abbiamo ritenuto opportuno (per evitare di preoccupare troppo il conducente) considerare un numero di secondi di massimo due cifre. Prima di terminare la funzione si controlla che il valore di \texttt{numb} non sia superiore a 99, in caso si salta a \texttt{\_numb\_overflow} che azzera \texttt{numb}.
	\end{itemize}
	
	\begin{framescelte}{Scelte Progettuali}
		Si è scelto di dividere il progetto in poche funzioni e di tenere tutti i controlli relativi a input/output in un'unica funzione chiamata \texttt{\_check}. Sarebbe stato tranquillamente possibile fare una divisione in funzioni minori per aumentare la leggibilità (e.g. \texttt{check\_rpm}, \texttt{check\_init}, \ldots), ma è stato ritenuto più pratico, per la semplicità dei controlli, mantenere il codice in un unico file. 
	\end{framescelte}
	
	\section{write\_line.s}
	Contiene la funzione che si occupa di creare la stringa di output e scriverla sul corrispondente file. 
	Ricordiamo qui la codifica utilizzata delle modalità di funzionamento e la struttura di una singola riga di output:
	\begin{table}[!h]
		\centering
		\begin{tabular}{| l | c | c |}
			\hline
			Spento & 0 & 00 \\ \hline
			SG & 1 & 01 \\ \hline
			OPT & 2 & 10 \\ \hline
			FG & 3 & 11 \\ \hline
		\end{tabular}
		\\
		\vspace{5mm}
		\begin{tabular}{| c | c | c | c | c | c |}
			\hline
			\texttt{alm} (1 byte) & , & \texttt{mod} (2 byte) & , & \texttt{numb} (2 byte) &  \texttt{\textbackslash n} \\ \hline
		\end{tabular}
	\end{table}
	\subsection{Variabili Locali}
	\begin{itemize}
		\itemtt{\hyperref[v:6:1]{output\_buff}}: Buffer di dimensione \texttt{OUTPUT\_BUFF\_LEN} che conterrà i caratteri della stringa da scrivere sul file di output.
		\itemtt{\hyperref[v:6:2]{OUTPUT\_BUFF\_LEN}}: Dimensione del buffer.
		\itemtt{\hyperref[v:6:3]{MOD\_XX}}: Stringhe costanti che identificano la modalità di funzionamento corrispondente in binario.
		\itemtt{\hyperref[v:6:4]{MOD\_LEN}}: Dimensione delle stringhe di modalità.
	\end{itemize}
	
	\subsection{Funzioni ed Etichette}
	\begin{itemize}
		\itemtt{\hyperref[e:6:1]{\_write\_line}}: Inizializza i registri necessari alla scrittura sulla variabile buffer.
		\itemtt{\hyperref[e:6:2]{\_alm\_X}}: Aggiunge al buffer il corretto valore di \texttt{alm}.
		\itemtt{\hyperref[e:6:3]{\_print\_mod}}: Aggiunge al buffer una virgola come separatore e in base al valore di \texttt{mod} salta alla corrispondente etichetta.
		\itemtt{\hyperref[e:6:4]{\_mod\_X}}: La stringa corrispondente alla modalità X codificata in binario viene messa in \texttt{eax}, in seguito viene aggiunta al buffer nell'etichetta \texttt{\_end\_print\_mod}.
		\itemtt{\hyperref[e:6:5]{\_print\_numb}}: Si occupa di aggiungere al buffer il valore di \texttt{numb} opportunamente convertito in \texttt{ASCII} e di terminare la stringa con il carattere \texttt{\textbackslash n}.
		\itemtt{\hyperref[e:6:6]{\_numb\_one\_digit}}: Se il valore di \texttt{numb} è minore di 10 si occupa di aggiungere uno 0 come prima cifra.
	\end{itemize}
	
	\section{itoa.s}
	Contiene la funzione per convertire un valore da intero a una corrispondente stringa \texttt{ASCII}.
	\subsection{Funzioni ed Etichette} 
	\begin{itemize}
		\itemtt{\hyperref[e:7:1]{\_itoa}}: Inizializza i registri necessari alla conversione.
		\itemtt{\hyperref[e:7:2]{\_itoa\_dividi}}: Si occupa di contare i caratteri necessari per la stringa e di posizionare il carattere \texttt{\textbackslash 0} alla fine della stringa.
		\itemtt{\hyperref[e:7:3]{\_itoa\_converti}}: Scrive ogni cifra nella posizione corretta della stringa.
	\end{itemize}
	
	\section{close\_files.s}
	Contiene la funzione per chiudere correttamente un file. 
	
	
	\newpage
	
	
		\section{Codice Sorgente}
		
		\subsection{syscall.inc}
		
		\begin{lstlisting}[language=MyAssembler, style=MyAsm]
		# file di definizione delle chiamate di sistema linux
			
		.equ    SYS_EXIT, 1(*@\label{s:1}@*)
		.equ    SYS_READ, 3
		.equ    SYS_WRITE, 4
		.equ    SYS_OPEN, 5
		.equ    SYS_CLOSE, 6
			
		.equ    STDIN, 0
		.equ    STDOUT, 1
		.equ    STDERR, 2
			
		.equ    SYSCALL, 0x80
		\end{lstlisting}
		\subsection{main.s}
		%	\lstinputlisting[style=MyAsm, language = MyAssembler]{../src/main.s}
		\begin{lstlisting}[language=MyAssembler, style=MyAsm]
		# Progetto Assembly 2016
		# File: main.s
		# Autori: Noè Murr, Mirko Morati
		#
		# Descrizione: File principale, punto di inizio del programma.
		.include    "syscall.inc"
		
		.section    .data
		input_fd:   .long 0(*@\label{v:1:1}@*)         # variabile globale che conterrà il file
		# descriptor del file di input
		
		output_fd:  .long 0(*@\label{v:1:2}@*)         # variabile globale che conterrà il file
		# descriptor del file di output
		
		# Variabili globali per i segnali di input
		init:   .long 0(*@\label{v:1:3}@*)
		reset:  .long 0(*@\label{v:1:4}@*)
		rpm:    .long 0(*@\label{v:1:5}@*)
		
		# Variabili globali per i segnali di output
		alm:    .long 0(*@\label{v:1:6}@*)
		numb:   .long 0(*@\label{v:1:8}@*)
		mod:    .long 0(*@\label{v:1:7}@*)
		
		# Codice del programma
		
		.section    .text
		.globl  input_fd
		.globl  output_fd
		.globl  init
		.globl  reset
		.globl  rpm
		.globl  alm
		.globl  numb
		.globl  mod
		.globl  _start
		
		# Stringa per mostrare l'utilizzo del programma in caso di parametri errati
		usage(*@\label{v:1:9}@*):  .asciz "usage: programName inputFilePath outputFilePath\n"
		.equ    USAGE_LENGTH(*@\label{v:1:10}@*), .-usage
		
		_start:(*@\label{e:1:1}@*)
		# Recupero i parametri del main
		popl    %eax                # Numero parametri
		
		# Controllo argomenti, se sbagliati mostro l'utilizzo corretto
		cmpl    $3, %eax
		jne     _show_usage
		
		popl    %eax                # Nome programma
		popl    %eax                # Primo parametro (nome file di input)
		popl    %ebx                # Secondo parametro (nome file di output)
		
		# NB: non salvo ebp in quanto non ha alcuna utilità
		# nella funzione start che comunque non ritorna
		
		movl    %esp, %ebp
		
		call    _open_files         # Apertura dei file
		
		_main_loop:(*@\label{e:1:2}@*)
		
		call    _read_line          # Leggiamo la riga
		
		cmpl    $-1, %ebx           # EOF se ebx == -1
		je      _end
		
		call    _check              # Controllo delle variabili
		
		call    _write_line         # Scrittura delle variabili di output su file
		
		jmp     _main_loop          # Leggi un altra riga finché non è EOF
		
		_end:(*@\label{e:1:3}@*)
		
		call    _close_files        # Chiudi correttamente i file
		
		# sys_exit(0);
		movl    $SYS_EXIT, %eax
		movl    $0, %ebx
		int     $SYSCALL
		
		_show_usage:(*@\label{e:1:4}@*)
		# esce in caso di errore con codice 1
		# sys_write(stdout, usage, USAGE_LENGTH);
		movl    $SYS_WRITE, %eax
		movl    $STDOUT, %ebx
		movl    $usage, %ecx
		movl    $USAGE_LENGTH, %edx
		int     $SYSCALL
		
		# sys_exit(1);
		movl    $SYS_EXIT, %eax
		movl    $1, %ebx
		int     $SYSCALL
		
		\end{lstlisting}
		
		\subsection{open\_files.s}
		\begin{lstlisting}[language=MyAssembler, style=MyAsm]
		# Progetto Assembly 2016
		# File: open_files.s
		# Autori: Noè Murr, Mirko Morati
		#
		# Descrizione: File contenente la funzione che si occupa di aprire i file di input e di
		# output, i file descriptor vengono inseriti in variabili globali.
		# Si suppone che il nome dei due file siano salvati negli indirizzi contenuti
		# rispettivamente in %eax (input) ed in %ebx (output).
		
		.include "syscall.inc"
		
		.section .text
		
		error_opening_files:(*@\label{v:2:1}@*) .asciz "errore nell' apertura dei file\n"
		.equ	ERROR_OPENING_LENGTH(*@\label{v:2:2}@*), .-error_opening_files
		
		.globl	_open_files				# Dichiaro la funzione globale
		.type 	_open_files, @function 	# Dichiaro l'etichetta come una funzione
		
		_open_files:(*@\label{e:2:1}@*)
		
		pushl	%ebp
		movl 	%esp, %ebp
		
		pushl	%ebx					# Pusho l' indirizzo del file di output
		# sullo stack
		
		movl 	%eax, %ebx 				# Sposto l' indirizzo del file che vado
		# ad aprire in %ebx
		
		movl 	$SYS_OPEN, %eax 		# Chiamata di sistema open
		movl 	$0, %ecx 				# read-only mode
		int 	$SYSCALL 				# Apro il file
		
		cmpl 	$0, %eax
		jl		_error_opening_files
		
		movl 	%eax, input_fd			# Metto il file descriptor nella
		# sua variabile
		
		popl 	%ebx 					# Riprendo l' indirizzo del nome del file
		# di output che avevo messo sullo stack
		
		movl 	$SYS_OPEN, %eax 		# Chiamata di sistema open
		movl 	$01101, %ecx 			# read and write mode
		movl    $0666, %edx            	# flags
		int 	$SYSCALL 				# Apro il file
		
		cmpl 	$0, %eax
		jl		_error_opening_files
		
		movl 	%eax, output_fd			# Metto il file descriptor nella
		# sua variabile
		
		movl 	%ebp, %esp
		popl	%ebp
		ret 							# Ritorna al chiamante
		
		_error_opening_files:(*@\label{e:2:2}@*)
		# Esce con codice di errore 2
		# sys_write(stdout, usage, USAGE_LENGTH);
		movl 	$SYS_WRITE, %eax
		movl 	$STDERR, %ebx
		movl 	$error_opening_files, %ecx
		movl 	$ERROR_OPENING_LENGTH, %edx
		int 	$SYSCALL
		
		# sys_exit(2);
		movl 	$SYS_EXIT, %eax
		movl 	$2, %ebx
		int 	$SYSCALL
		\end{lstlisting}
		
		\subsection{read\_line.s}
		\begin{lstlisting}[language=MyAssembler, style=MyAsm]
		# Progetto Assembly 2016
		# File: read_line.s
		# Autori: Noè Murr, Mirko Morati
		#
		# Descrizione: Funzione che legge una riga alla volta del file di input.
		
		.include "syscall.inc"
		
		.section .bss
		.equ	INPUT_BUFF_LEN, 9(*@\label{v:3:2}@*)
		input_buff:	.space INPUT_BUFF_LEN(*@\label{v:3:1}@*)	# Input buffer di 9 byte
		
		.section .text
		.globl 	_read_line
		.type   _read_line, @function
		
		_read_line:(*@\label{e:3:1}@*)
		pushl	%ebp
		movl 	%esp, %ebp
		
		# Lettura riga
		# sys_read(input_fd, input_buff, INPUT_BUFF_LEN);
		movl    input_fd, %ebx
		movl    $SYS_READ, %eax
		leal    input_buff, %ecx
		movl    $INPUT_BUFF_LEN, %edx
		int     $SYSCALL
		
		cmpl    $0, %eax                    # Se eax == 0 EOF
		je      _eof
		
		# Estrazione dei valori di init, reset, rpm dal buffer
		leal    input_buff, %edi
		call    _atoi
		movl    %eax, init
		
		incl    %edi						# Salto il carattere ','
		
		call    _atoi
		movl    %eax, reset
		
		incl    %edi						# Salto il carattere ','
		
		call    _atoi
		movl    %eax, rpm
		
		movl    %ebp, %esp
		popl    %ebp
		
		xorl    %ebx, %ebx					# ebx = 0 permette di proseguire
		ret
		
		_eof:(*@\label{e:3:2}@*)
		# in caso di EOF %ebx = -1
		movl    %ebp, %esp
		popl    %ebp
		
		movl    $-1, %ebx
		ret
		\end{lstlisting}
		
		\newpage
		\subsection{atoi.s}
		\begin{lstlisting}[language=MyAssembler, style=MyAsm]
		# Progetto Assembly 2016
		# File: atoi.s
		# Autori: Alessandro Righi, Noè Murr, Mirko Morati
		#
		# Descrizione: Funzione che converte una stringa in intero.
		
		.section .text
		.globl	_atoi
		.type	_atoi, @function
		
		# Funzione che converte una stringa di input in numero
		# Prototipo C-style:
		#   uint32_t atoi(const char *string);
		# Parametri di input:
		#   EDI - Stringa da convertire
		# Parametri di output:
		#   EAX - Valore convertito
		
		_atoi:(*@\label{e:4:1}@*)
		xorl	%eax, %eax   # azzero il registro EAX per contenere il risultato
		xorl	%ebx, %ebx   # azzero EBX
		movl	$10, %ecx    # sposto 10 in ECX che conterrà il valore moltiplicativo
		
		_atoi_loop:(*@\label{e:4:2}@*)
		xorl	%ebx, %ebx
		movb	(%edi), %bl  # sposto un byte dalla stringa in BL
		subb	$48, %bl     # sottraggo il valore ASCII dello 0 a BL
		# per avere un valore intero
		
		cmpb	$0, %bl      # Se il numero è minore di 0
		jl		_atoi_end    # allora esco dal ciclo
		cmpb	$10, %bl     # Se il numero è maggiore o uguale a 10
		jge 	_atoi_end    # esco dal ciclo
		
		mull	%ecx         # altrimenti moltiplico EAX per 10
		# (10 messo precedentemente in ECX)
		addl	%ebx, %eax   # aggiungo a EAX il valore attuale
		incl	%edi         # incremento EDI
		
		jmp		_atoi_loop   # rieseguo il ciclo
		
		_atoi_end:
		ret
		\end{lstlisting}
		
		\newpage
		\subsection{check.s}
		\begin{lstlisting}[language=MyAssembler, style=MyAsm]
		# Progetto Assembly 2016
		# File: check.s
		# Autori: Noè Murr, Mirko Morati
		#
		# Descrizione: Funzione che controlla le variabili
		# init, reset, rpm e setta le variabili alm, mod e numb
		
		.section .data
		
		.section .text
		.globl  _check
		.type   _check, @function
		
		_check:(*@\label{e:5:1}@*)
		pushl   %ebp
		movl    %esp, %ebp
		
		# Caso init == 0: alm = 0; mod = 0; numb = 0;
		cmpl    $0, init
		je      _init_0
		
		# Caso SG: alm = 0; mod = 1; numb = reset == 1 ? 0 : numb + 1;
		cmpl    $2000, rpm
		jl      _sg
		
		# Caso OPT: alm = 0; mod = 2; numb = reset == 1 ? 0 : numb + 1;
		cmpl    $4000, rpm
		jle     _opt
		
		# Caso FG: alm = numb >= 15? 1 : 0; mod = 3; numb = reset == 1 ? 0 : numb + 1;
		_fg:(*@\label{e:5:2}@*)
		# Salviamo la nuova modalita' in %eax e controlliamo reset
		movl    $3, %eax
		cmpl    $1, reset
		je      _reset_numb
		
		# Se la nuova modalità non è la stessa si resetta il numero di secondi
		cmpl    $3, mod
		jne     _reset_numb
		
		incl    numb
		movl    %eax, mod
		
		# Se il numero di secondi è maggiore o uguale a 15 viene alzata l'allarme
		cmpl    $15, numb
		jge      _set_alm
		
		jmp     _end_check
		
		_opt:
		movl    $2, %eax
		cmpl    $1, reset
		je      _reset_numb
		
		cmpl    $2, mod
		jne     _reset_numb
		
		incl    numb
		movl    %eax, mod
		
		jmp     _end_check
		
		_sg:
		movl    $1, %eax
		cmpl    $1, reset
		je      _reset_numb
		
		cmpl    $1, mod
		jne      _reset_numb
		
		incl    numb
		movl    %eax, mod
		
		jmp     _end_check
		
		_reset_numb:(*@\label{e:5:3}@*)
		movl    %eax, mod
		movl    $0, numb
		movl    $0, alm
		
		jmp     _end_check
		
		_set_alm:(*@\label{e:5:4}@*)
		movl    $1, alm
		
		
		jmp     _end_check
		
		_init_0:(*@\label{e:5:5}@*)
		movl    $0, alm
		movl    $0, numb
		movl    $0, mod
		
		_end_check:(*@\label{e:5:6}@*)
		
		# Se il numero di secondi supera i 99 allora dobbiamo ricominciare il conteggio
		cmpl    $99, numb
		jg      _numb_overflow
		movl    %ebp, %esp
		popl    %ebp
		
		ret
		
		_numb_overflow:(*@\label{e:5:7}@*)
		movl    $0, numb
		jmp     _end_check
		\end{lstlisting}
		
		\subsection{write\_line.s}
		\begin{lstlisting}[language=MyAssembler, style=MyAsm]
		# Progetto Assembly 2016
		# File: check.s
		# Autori: Noè Murr, Mirko Morati
		#
		# Descrizione: Funzione che scrive una riga alla volta nel file di output
		
		.include "syscall.inc"
		
		.section .bss
		.equ	OUTPUT_BUFF_LEN, 8(*@\label{v:6:2}@*)
		output_buff:(*@\label{v:6:1}@*) .space OUTPUT_BUFF_LEN
		
		.section .text
		.globl  _write_line
		.type   _write_line, @function
		MOD_00: .ascii "00"(*@\label{v:6:3}@*)           # motore spento
		MOD_01: .ascii "01"           # motore sotto giri
		MOD_10: .ascii "10"           # motore in stato ottimale
		MOD_11: .ascii "11"           # motore fuori giri
		.equ 	MOD_LEN, 2(*@\label{v:6:4}@*)
		
		_write_line:(*@\label{e:6:1}@*)
		pushl   %ebp
		movl    %esp, %ebp
		
		leal    output_buff, %edi     # spostiamo il puntatore
		# del buffer di output in EDI
		
		cmpl    $1, alm               # se l'allarme è 1 stampiamo 1
		# altrimenti 0 senza chiamare funzioni
		je      _alm_1
		
		_alm_0:(*@\label{e:6:2}@*)
		movl    $48, (%edi)
		jmp     _print_mod
		
		_alm_1:
		movl    $49, (%edi)
		
		_print_mod:(*@\label{e:6:3}@*)
		movl    $44, 1(%edi)          # aggiungiamo la virgola dopo
		# il segnale di allarme
		addl    $2, %edi              # spostiamo un immaginario cursore
		# nella posizione dove stampare la mod
		
		cmpl    $1, mod               # controlliamo il valore di mod
		# e stampiamo la stringa corretta in base
		# alla giusta modalita' di funzionamento
		je      _mod_1
		cmpl    $2, mod
		je      _mod_2
		cmpl    $3, mod
		je      _mod_3
		
		_mod_0:(*@\label{e:6:4}@*)
		movl    MOD_00, %eax
		jmp _end_print_mod
		
		_mod_1:
		movl    MOD_01, %eax
		jmp _end_print_mod
		
		_mod_2:
		movl    MOD_10, %eax
		jmp _end_print_mod
		
		_mod_3:
		movl    MOD_11, %eax
		
		_end_print_mod:
		movl    %eax, (%edi)          # mettiamo la stringa nell' output_buff
		addl    $MOD_LEN, %edi        # spostato il cursore (la posizione di edi)
		# nel punto esatto dove scrivere
		movl    $44, (%edi)           # aggiungiamo la virgola
		incl    %edi                  # spostiamo il cursore
		
		cmpl    $10, numb             # controlliamo se il numero di secondi
		# è ad una sola cifra, in tal caso
		# aggiungiamola cifra 0
		jl      _numb_one_digit
		
		_print_numb:(*@\label{e:6:5}@*)
		movl    numb, %eax            # prepariamo la chiamata per itoa
		
		call    _itoa                 # chiamiamo itoa
		
		
		leal    output_buff, %edi     # mettiamo il puntatore di output_buff in edi
		addl    $7, %edi              # ci aggiungiamo 7 per arrivare
		# alla fine della stringa,
		movl    $10, (%edi)           # punto nel quale aggiungiamo un \n
		
		movl    $SYS_WRITE, %eax
		movl    output_fd, %ebx
		leal    output_buff, %ecx
		movl    $OUTPUT_BUFF_LEN, %edx
		int     $SYSCALL
		
		
		movl    %ebp, %esp
		popl    %ebp
		
		ret
		
		_numb_one_digit:(*@\label{e:6:6}@*)
		movl    $48, (%edi)
		incl    %edi
		jmp     _print_numb
		\end{lstlisting}
		
		\subsection{itoa.s}
		\begin{lstlisting}[language=MyAssembler, style=MyAsm]
		# Progetto Assembly 2016
		# File: itoa.s
		# Autori: Alessandro Righi, Noè Murr, Mirko Morati
		#
		# Descrizione: Funzione che converte un intero in stringa
		# Prototipo C-style:
		#   u_int32_t itoa(uint32_t val, char *string);
		# Parametri di input:
		#   EAX - Valore intero unsigned a 64bit da convertire
		#   EDI - Puntatore alla stringa su cui salvare il risultato
		# Parametri di output:
		#   EAX - Lunghezza della stringa convertita (compresiva di \0 finale)
		
		.section .text
		.global _itoa
		.type   _itoa, @function
		
		_itoa:(*@\label{e:7:1}@*)
		movl    $10, %ecx   # porto il fattore moltiplicativo in ECX
		movl    %eax, %ebx  # salvo temporaneamente il valore di EAX in EBX
		xorl    %esi, %esi  # azzero il registro ESI
		
		_itoa_dividi:(*@\label{e:7:2}@*)
		xorl    %edx, %edx  # azzero EDX per fare la divisione
		divl    %ecx        # divide EAX per ECX, salva il resto in EDX
		incl    %esi        # incrementa il contatore
		testl   %eax, %eax  # se il valore di EAX non è zero ripeti il ciclo
		jnz     _itoa_dividi
		
		addl    %esi, %edi  # somma all'indirizzo del buffer
		# il numero di caratteri del numero
		movl    %ebx, %eax  # rimette il valore da convertire in EAX
		movl    %esi, %ebx  # salvo il valore della lunghezza della stringa in EBX
		
		movl    $0, (%edi)  # aggiungo un null terminator alla fine della stringa
		decl    %edi        # decremento il contatore della stringa di 1
		
		_itoa_converti:(*@\label{e:7:3}@*)
		xorl    %edx, %edx  # azzero EDX per fare la divisione
		divl    %ecx        # divido EAX per ECX, salvo il valore del resto in EDX
		addl    $48, %edx   # sommo 48 a EDX
		movb    %dl, (%edi) # sposto il byte inferiore di EDX (DL)
		# nella locazione di memoria puntata da EDI
		decl    %edi        # decremento il puntatore della stringa
		decl    %esi        # decremento il contatore
		testl   %esi, %esi  # se il contatore non è 0 continua ad eseguire il loop
		jnz     _itoa_converti
		
		movl    %ebx, %eax  # porto il valore della lunghezza
		# della stringa in EAX per ritornarlo
		incl    %eax        # incremento di 1 EAX (in modo da includere il \0)
		ret
		\end{lstlisting}
		
		\subsection{close\_files.s}
		\begin{lstlisting}[language=MyAssembler, style=MyAsm]
		# Progetto Assembly 2016
		# File: check.s
		# Autori: Noè Murr, Mirko Morati
		#
		# Descrizione: Funzione che chiude i file aperti precedentemente
		
		.include "syscall.inc"
		
		.section .text
		.globl 	_close_files 			 # Dichiaro la funzione globale
		.type 	_close_files, @function  # Dichiaro l' etichetta come una funzione
		
		_close_files:
		
		pushl	%ebp
		movl 	%esp, %ebp
		
		# sys_close(input_fd);
		movl	$SYS_CLOSE, %eax
		movl	input_fd, %ebx
		int 	$SYSCALL
		
		# sys_close(output_fd);
		movl	$SYS_CLOSE, %eax
		movl	output_fd, %ebx
		int 	$SYSCALL
		
		
		movl    %ebp, %esp
		popl    %ebp
		
		ret
		\end{lstlisting}
		
	\end{document}